% GENERATED FILE. Edit canonical lessons, then run npm run generate:books.
% canonical-source-hash: fnv1a64:cfc3646b6feb2cd5
% canonical-lessons: MR-C82-phekne, MR-C82-udne, MR-C82-pohne, MR-C82-shikavne, MR-C82-banavne

\chapter{To throw, To fly, To swim, To teach, To make}
\label{ch:mr-to-throw-to-fly-to-swim-to-teach-to-make}
\hlchaptermodality{82}

\begin{chapteropening}
\textbf{By the end of this chapter:} \emph{I can say to throw, to fly, to swim, to teach and to make.}

Complete the last lesson of chapter 82: I can say to throw, to fly, to swim, to teach and to make.
\end{chapteropening}

\section[\texorpdfstring{\mr{फेकणे} (phekṇe)}{फेकणे (phekṇe)}]{\mr{फेकणे} (phekṇe) — to throw}

\label{lesson:MR-C82-phekne}



\begin{quote}
Before the new one: say the Marathi for to put on, then the Marathi for to leave.
\end{quote}

\subsection*{You'll want to know: \mr{फेकणे}}
\textbf{\mr{फेकणे}} — \emph{phekṇe} — \textquotedblleft{}to throw\textquotedblright{}.

\begin{cousinweb}
\textbf{\mr{फेकणे}} (\emph{phekṇe}) comes from \emph{*phēṅk-}, \textquotedblleft{}throw\textquotedblright{}, a form CDIAL reconstructs beside \emph{*phēkk-} (CDIAL 9106).
\end{cousinweb}

\begin{grammarlens}[title={What you've built}]
One more word for this chapter.
\end{grammarlens}

\subsection*{Your turn}
\begin{itemize}
\raggedright
  \item \emph{Say it:} \emph{phekṇe}
  \item \emph{Say it:} \emph{phekṇe}, once more
  \item \emph{Say it:} say \emph{phekṇe}
  \item \emph{Recall:} say the Marathi for to put on, then the Marathi for to leave, then say \emph{phekṇe} again
\end{itemize}

\begin{tcolorbox}[breakable,colback=teal!4,colframe=teal!35!black,title={Before you move on}]
What does \mr{फेकणे} mean? (\textquotedblleft{}To throw\textquotedblright{}.) Say it once more.
\end{tcolorbox}
\section[\texorpdfstring{\mr{उडणे} (uḍṇe)}{उडणे (uḍṇe)}]{\mr{उडणे} (uḍṇe) — to fly}

\label{lesson:MR-C82-udne}



\begin{quote}
Before the new one: say the Marathi for to leave, then the Marathi for to throw.
\end{quote}

\subsection*{You'll want to know: \mr{उडणे}}
\textbf{\mr{उडणे}} — \emph{uḍṇe} — \textquotedblleft{}to fly\textquotedblright{}.

\begin{cousinweb}
\textbf{\mr{उडणे}} (\emph{uḍṇe}) is a tadbhava, a word inherited from Sanskrit and reshaped by sound change. It comes from \emph{uḍḍayate}, \textquotedblleft{}flies up\textquotedblright{} (CDIAL 1697).
\end{cousinweb}

\begin{grammarlens}[title={What you've built}]
One more word for this chapter.
\end{grammarlens}

\subsection*{Your turn}
\begin{itemize}
\raggedright
  \item \emph{Say it:} \emph{uḍṇe}
  \item \emph{Say it:} \emph{uḍṇe}, once more
  \item \emph{Say it:} say \emph{uḍṇe}
  \item \emph{Recall:} say the Marathi for to leave, then the Marathi for to throw, then say \emph{uḍṇe} again
\end{itemize}

\begin{tcolorbox}[breakable,colback=teal!4,colframe=teal!35!black,title={Before you move on}]
What does \mr{उडणे} mean? (\textquotedblleft{}To fly\textquotedblright{}.) Say it once more.
\end{tcolorbox}
\section[\texorpdfstring{\mr{पोहणे} (pohṇe)}{पोहणे (pohṇe)}]{\mr{पोहणे} (pohṇe) — to swim}

\label{lesson:MR-C82-pohne}



\begin{quote}
Before the new one: say the Marathi for to throw, then the Marathi for to fly.
\end{quote}

\subsection*{You'll want to know: \mr{पोहणे}}
\textbf{\mr{पोहणे}} — \emph{pohṇe} — \textquotedblleft{}to swim\textquotedblright{}.

\begin{cousinweb}
\textbf{\mr{पोहणे}} (\emph{pohṇe}) is a tadbhava, a word inherited from Sanskrit and reshaped by sound change. It comes from \emph{pravahati}, \textquotedblleft{}carries forward; drives on\textquotedblright{} (CDIAL 8792). CDIAL glosses the Marathi verb \textquotedblleft{}to swim, float\textquotedblright{}.
\end{cousinweb}

\begin{grammarlens}[title={What you've built}]
One more word for this chapter.
\end{grammarlens}

\subsection*{Your turn}
\begin{itemize}
\raggedright
  \item \emph{Say it:} \emph{pohṇe}
  \item \emph{Say it:} \emph{pohṇe}, once more
  \item \emph{Say it:} say \emph{pohṇe}
  \item \emph{Recall:} say the Marathi for to throw, then the Marathi for to fly, then say \emph{pohṇe} again
\end{itemize}

\begin{tcolorbox}[breakable,colback=teal!4,colframe=teal!35!black,title={Before you move on}]
What does \mr{पोहणे} mean? (\textquotedblleft{}To swim\textquotedblright{}.) Say it once more.
\end{tcolorbox}
\section[\texorpdfstring{\mr{शिकवणे} (shikavṇe)}{शिकवणे (shikavṇe)}]{\mr{शिकवणे} (shikavṇe) — to teach}

\label{lesson:MR-C82-shikavne}



\begin{quote}
Before the new one: say the Marathi for to fly, then the Marathi for to swim.
\end{quote}

\subsection*{You'll want to know: \mr{शिकवणे}}
\textbf{\mr{शिकवणे}} — \emph{shikavṇe} — \textquotedblleft{}to teach\textquotedblright{}.

\begin{cousinweb}
\textbf{\mr{शिकवणे}} (\emph{shikavṇe}) is the \textquotedblleft{}make\textquotedblright{} form of \textbf{\mr{शिकणे}}, \textquotedblleft{}to learn\textquotedblright{}: teaching is making someone learn. CDIAL lists both under Sanskrit \emph{śikṣate}, \textquotedblleft{}learns\textquotedblright{} (CDIAL 12430).
\end{cousinweb}

\begin{grammarlens}[title={What you've built}]
One more word for this chapter.
\end{grammarlens}

\subsection*{Your turn}
\begin{itemize}
\raggedright
  \item \emph{Say it:} \emph{shikavṇe}
  \item \emph{Say it:} \emph{shikavṇe}, once more
  \item \emph{Say it:} say \emph{shikavṇe}
  \item \emph{Recall:} say the Marathi for to fly, then the Marathi for to swim, then say \emph{shikavṇe} again
\end{itemize}

\begin{tcolorbox}[breakable,colback=teal!4,colframe=teal!35!black,title={Before you move on}]
What does \mr{शिकवणे} mean? (\textquotedblleft{}To teach\textquotedblright{}.) Say it once more.
\end{tcolorbox}
\section[\texorpdfstring{\mr{बनवणे} (banavṇe)}{बनवणे (banavṇe)}]{\mr{बनवणे} (banavṇe) — to make}

\label{lesson:MR-C82-banavne}



\begin{quote}
Before the new one: say the Marathi for to swim, then the Marathi for to teach.
\end{quote}

\subsection*{You'll want to know: \mr{बनवणे}}
\textbf{\mr{बनवणे}} — \emph{banavṇe} — \textquotedblleft{}to make\textquotedblright{}.

\begin{cousinweb}
\textbf{\mr{बनवणे}} (\emph{banavṇe}) is the \textquotedblleft{}make\textquotedblright{} form of Marathi \emph{banṇẽ}, \textquotedblleft{}to suit, be made\textquotedblright{}. CDIAL lists that verb under Sanskrit \emph{vanati}, \textquotedblleft{}desires, makes ready\textquotedblright{}, and marks it as taken from Hindi; Bloch's source, \emph{varṇayati}, it says does not explain the word's sounds (CDIAL 11260).
\end{cousinweb}

\begin{grammarlens}[title={What you've built}]
One more word for this chapter.
\end{grammarlens}

\subsection*{Your turn}
\begin{itemize}
\raggedright
  \item \emph{Say it:} \emph{banavṇe}
  \item \emph{Say it:} \emph{banavṇe}, once more
  \item \emph{Say it:} say \emph{banavṇe}
  \item \emph{Recall:} say the Marathi for to swim, then the Marathi for to teach, then say \emph{banavṇe} again
\end{itemize}

\begin{tcolorbox}[breakable,colback=teal!4,colframe=teal!35!black,title={Before you move on}]
What does \mr{बनवणे} mean? (\textquotedblleft{}To make\textquotedblright{}.) Say it once more.
\end{tcolorbox}
